\chapter*{Abstract}

In these lecture notes, we consider the foundational principles of Quantum Information Theory, the
interdisciplinary framework governing the storage, processing, and transmission of information via
quantum systems. The material is pedagogically divided into two conceptually distinct regimes. In
Part I, we establish the theoretical baseline of idealized, isolated quantum systems, focusing on
the paradigm shift from classical Boolean logic to the reversible, unitary evolution of quantum
states. Within this framework, we formalize the mathematical description of the quantum bits, or
qubits, superposition and entanglement within finite-dimensional Hilbert spaces. This serves as the
mathematical basis for defining fundamental quantum circuits, communication protocols and canonical
algorithms. In Part II, we extend the formalism beyond perfectly isolated systems, detailing the
transition from pure mathematical abstraction to the physical reality of open quantum systems. We
employ density matrices to model environmental coupling, quantum noise, and decoherence, and we
analyze the von Neumann entropy and the thermodynamics of information. Because of the inherent
fragility of these physical states, we focus on the necessity of quantum error-correction protocols.
Finally, we map this theoretical formalism to active experimental architectures defining the noisy
intermediate-scale quantum (NISQ) era, analyzing the physical realization of quantum computers using
superconducting circuits, trapped ions, and topological qubits.
